% !TeX program = xelatex
\documentclass[11pt,a4paper]{article}
\usepackage[margin=20mm]{geometry}
\usepackage{fontspec}
\IfFontExistsTF{Times New Roman}{\setmainfont{Times New Roman}}{%
\setmainfont{NimbusRoman-Regular.otf}[BoldFont=NimbusRoman-Bold.otf,
ItalicFont=NimbusRoman-Italic.otf,BoldItalicFont=NimbusRoman-BoldItalic.otf]}
\usepackage{setspace}\setstretch{1.08}
\usepackage{titlesec}
\titlespacing*{\subsection}{0pt}{9pt}{4pt}
\usepackage{array,tabularx,longtable,booktabs,enumitem}
\usepackage[hidelinks]{hyperref}\usepackage{xurl}\urlstyle{same}
\setlength{\parindent}{0pt}\setlength{\parskip}{4pt}
\setlist{leftmargin=5mm,topsep=3pt,itemsep=1pt,parsep=0pt}
\newcolumntype{Y}{>{\raggedright\arraybackslash}X}
\newcommand{\Page}[1]{\clearpage\section*{#1}}
\newcommand{\Key}[1]{\textbf{#1}}
\begin{document}
{\Large\bfseries BE Project Report Preparation Instructions}\par
VESIT\quad Department of Automation and Robotics\par
Student handout\quad 23 October 2026

Prepare a report that explains your own engineering work and provides verifiable
evidence for its results. Follow the supplied report-format PDF for institutional
requirements. The chapter plan and evidence checklists below are recommendations;
your guide should approve adaptations.

\subsection*{1 Start the LaTeX project}
\begin{enumerate}
\item Upload the ZIP to Overleaf as a new project, or extract it for local use.
Set \texttt{main.tex} as the main document and select \Key{XeLaTeX}.
\item Edit \texttt{project\_details.tex}. Enter the approved title, exact degree
and program wording, students and roll numbers, guide, department, institute and
submission year. Do not infer the award year from the academic year.
\item Place the approved institute emblem at
\texttt{figures/institute\_logo.png}. Do not reuse the other institute's
emblem in the supplied PDF specimen.
\item Replace all square-bracketed writing prompts. Four student placeholders
are provided; see the README for changing team size.
\item Compile until the contents and reference numbers settle. Locally run
\texttt{latexmk -xelatex main.tex}, or run \texttt{xelatex main.tex} at least twice.
\end{enumerate}

\subsection*{2 Required report format}
\begin{tabularx}{\textwidth}{@{}p{46mm}Y@{}}\toprule
Item & Requirement or selected implementation\\\midrule
Paper & A4, 210 by 297 mm; BE report printed on both sides\\
Body text & Times New Roman, 12 pt, 1.5 line spacing\\
Writing area & 160 mm wide by 245 mm high\\
Margins without header & Top 30 mm; bottom 22 mm; left 30 mm; right 20 mm\\
Page number & Centre of footer; Roman preliminary pages, Arabic main text\\
Headings & Chapter 18 pt centred; section 16 pt and subsection 14 pt bold, left aligned\\
Captions & Table title above; figure caption below; numbers by chapter\\
\bottomrule\end{tabularx}

\Key{Required font.} The included preview uses Nimbus Roman because Times New
Roman was unavailable during compilation. This substitute is not claimed to
meet the final font requirement. Install Times New Roman or upload legally
obtained \texttt{times.ttf}, \texttt{timesbd.ttf}, \texttt{timesi.ttf} and
\texttt{timesbi.ttf} into \texttt{fonts/}. The source selects Times New Roman
when it is available. Inspect the final PDF fonts before submission.

\Page{Order of Pages and Pagination}
\subsection*{3 Preliminary pages in order}
\begin{enumerate}
\item \Key{Inner title sheet}: first printed report page, following Specimen A.
Include the title, submission statement, degree, students and roll numbers,
supervisor, department, institute and submission year. Outer-cover artwork is separate.
\item \Key{Dedication}: optional. If included, it is auxiliary page i with its
printed number hidden. Disable it when not used.
\item \Key{Certificate}: internal approval sheet following Specimen B. With no
dedication, it is auxiliary page i, counted but unprinted. Obtain genuine signatures.
\item \Key{Project Report Approval for B.E.}: Specimen C, included only in the
post-viva final hardbound copy after successful examination and actual approval.
\item \Key{Declaration of Academic Honesty and Integrity}: Specimen D, adapted
for a team and signed by every student. It follows approval when included;
otherwise follows the certificate.
\item \Key{Abstract}: 500 words explaining the problem, method, important results
and contribution. Also provide a brief chapter-wise report layout and keywords.
\item \Key{Contents}: chapters, sections and subsections with page numbers at the right.
\item \Key{List of Figures}: figure numbers, captions and page numbers.
\item \Key{List of Tables}: table numbers, titles and page numbers.
\item \Key{Abbreviations, Notation and Nomenclature}: the complete list used in
this report, including Greek letters, subscripts, superscripts and relevant units.
\end{enumerate}

\subsection*{4 Main text and ending sections}
Introduction is the first chapter; Review of Literature is normally second.
Results and Discussions is the penultimate main chapter and Conclusions is last.
Then place \Key{Appendix I, Appendix II and so on}, followed by
\Key{Literature Cited}, optional \Key{Publications by the Students}, and
\Key{Acknowledgements} signed and dated by students. The appendix and
acknowledgement positions follow the supplied PDF.

\subsection*{5 Numbering and submission stage}
The title and optional outer cover have no printed page numbers. Preliminary
pages use continuous lower-case Roman numbering from the first auxiliary sheet,
with the first number hidden. Introduction resets to Arabic page 1, also hidden;
the next page shows 2. Subsequent pages continue sequentially. Chapters start on
fresh pages without automatically inserted blank recto pages.

Keep \texttt{\string\finalcopyfalse} before viva. Change it to
\texttt{\string\finalcopytrue} only for the final approved copy. Never pre-fill
examiner approval or signatures. The source contains a numbering typo in section
3.7; this template follows one continuous Roman sequence.

\Page{Suggested Contents for Chapters 1 to 4}
Here ``index'' means \Key{Table of Contents}. An alphabetical subject index is
not requested by the supplied guidelines. Adapt the technical headings to the
actual project with guide approval.

\subsection*{Chapter 1 Introduction}
\begin{itemize}
\item 1.1 Background and Motivation
\item 1.2 Problem Statement
\item 1.3 Aim and Objectives
\item 1.4 Scope and Limitations
\item 1.5 Project Contributions
\item 1.6 Organisation of the Report
\end{itemize}
Define the need, engineering problem and measurable objectives. Distinguish
original contributions from reused components or software.

\subsection*{Chapter 2 Review of Literature}
\begin{itemize}
\item 2.1 Review Method and Sources
\item 2.2 Related Work and Existing Systems
\item 2.3 Comparative Analysis
\item 2.4 Research or Engineering Gap
\end{itemize}
Critically compare relevant papers, standards, data sheets and products. Explain
the gap supported by the review rather than listing unrelated summaries.

\subsection*{Chapter 3 Problem Analysis and Requirements}
\begin{itemize}
\item 3.1 Stakeholders and Use Cases
\item 3.2 Functional Requirements
\item 3.3 Performance and Other Requirements
\item 3.4 Constraints and Feasibility
\item 3.5 Risks and Applicable Standards
\item 3.6 Requirements Traceability Plan
\end{itemize}
Assign requirement IDs and define acceptance criteria before validation.

\subsection*{Chapter 4 System Design and Methodology}
\begin{itemize}
\item 4.1 Proposed System Architecture
\item 4.2 Design Alternatives and Selection
\item 4.3 Hardware and Mechanical Design
\item 4.4 Software and Algorithm Design
\item 4.5 Mathematical Model and Design Calculations
\item 4.6 Safety Sustainability Ethics and Cost Decisions
\end{itemize}
Use calculations and comparisons to justify choices. Rename irrelevant headings,
for example hardware design in a software-only project.

\Page{Suggested Contents for Chapters 5 to 8 and Appendices}
\subsection*{Chapter 5 Implementation and Integration}
\begin{itemize}
\item 5.1 Development Environment and Tools
\item 5.2 Subsystem Implementation
\item 5.3 System Integration
\item 5.4 Prototype and Operating Procedure
\item 5.5 Implementation Challenges and Changes
\item 5.6 Bill of Materials and Actual Cost
\end{itemize}
Describe what was actually built and configured, with technical evidence,
software versions, implementation problems and corrective changes.

\subsection*{Chapter 6 Testing and Validation Methodology}
\begin{itemize}
\item 6.1 Test Plan and Acceptance Criteria
\item 6.2 Experimental Setup and Instruments
\item 6.3 Datasets and Test Conditions
\item 6.4 Test Procedures and Repeated Trials
\item 6.5 Performance Metrics and Analysis Method
\item 6.6 Fault Safety and Robustness Tests
\end{itemize}
Provide reproducible and safe procedures, trial counts, units, uncertainty,
data selection and operating conditions.

\subsection*{Chapter 7 Results and Discussions}
\begin{itemize}
\item 7.1 Measured Results
\item 7.2 Verification of Project Requirements
\item 7.3 Comparison with Baseline or Existing System
\item 7.4 Interpretation and Engineering Tradeoffs
\item 7.5 Safety Cost and Sustainability Evaluation
\item 7.6 Limitations and Contribution of the Work
\end{itemize}
Separate measured findings from estimates and simulations. Explain failed trials
and unmet requirements rather than omitting them.

\subsection*{Chapter 8 Conclusions}
8.1 Summary of the Work; 8.2 Conclusions against Objectives;
8.3 Scope for Future Work. Enumerate conclusions separately, supported by results.
Do not introduce new experimental claims in the conclusions.

\subsection*{Appendices and remaining entries}
\Key{Appendix I}: detailed design;
\Key{Appendix II}: raw tests and supporting data;
\Key{Appendix III}: software and reproducibility;
\Key{Appendix IV}: project management and supporting records.
Then Literature Cited, optional Publications, and Acknowledgements.
Meeting logs, individual contributions and outcome mapping are recommended
departmental evidence, not additional mandatory pages in the formatting PDF.

\Page{Writing in LaTeX and Maintaining References}
\subsection*{6 Generate contents automatically}
Use structural commands. Do not type contents page numbers by hand.
\begin{verbatim}
\chapter{System Design and Methodology}
\section{Proposed System Architecture}
\subsection{Sensor Interface}
\label{sec:sensor}
See Section~\ref{sec:sensor}.
\end{verbatim}
Add subsections where helpful; the template lists them as 4.1.1 and so on.
Deeper subdivisions need not be numbered or included in the contents.

\subsection*{7 Figures and tables}
Mention each figure or table before its first appearance. Keep small items near
that mention and use separate pages for large ones where appropriate. Label
axes and units. Cite reproduced material on a separate line in the caption.
\begin{verbatim}
\begin{figure}[htbp]
  \centering
  \includegraphics[width=.85\textwidth]{prototype.png}
  \caption{Prototype with labelled subsystems}
  \label{fig:prototype}
\end{figure}
\end{verbatim}
Place a table caption and label \Key{above} its tabular contents. Place a figure
caption and label \Key{below} its image. Use \texttt{longtable} for multipage
tables and repeat column headings. Do not use screenshots of tables or equations.

\subsection*{8 Citations and bibliography}
Use one consistent approved journal style, such as IEEE where appropriate. The
provided bibliography is a minimal editable scaffold, not real literature.
Replace its demonstration entry and citation.
\begin{verbatim}
Prior work is discussed in~\cite{motor-control}.
% In backmatter/references.tex:
\bibitem{motor-control}
... complete, verified bibliographic details ...
\end{verbatim}
Give all author names with initials, full title, source, volume, issue, pages and
year as applicable. Website citations need author/organisation, title, URL and
last-accessed date. Personal communications need author, title if any and receipt
date. Check every cited source; never invent bibliographic details.

\subsection*{9 Common source checks}
Escape special text characters: \verb|\&|, \verb|\%|, \verb|\_|,
\verb|\#| and \verb|\$|. Use \verb|\url{...}| for web links. Use numbered
equation environments and cross-references. Recompile after adding labels,
figures, citations or chapters. Resolve undefined references before submission.

\Page{Project Evidence and Final Submission Checklist}
\subsection*{10 Verification of project requirements}
Maintain the complete requirement-versus-result table in Results and Discussions.
For each requirement record:
\begin{enumerate}
\item Requirement ID and approved requirement.
\item Acceptance criterion, units and conditions.
\item Test ID, procedure, date, equipment and trial count.
\item Actual measured result; use Not tested if pending.
\item Status: Met, Partially met or Not met, with a reason.
\item Evidence reference to a test log, figure, appendix or versioned file.
\end{enumerate}
A vision project may need accuracy on an independent labelled dataset, throughput
under a defined load and robustness under defined lighting. Set justified targets
and report actual observations. Template examples must not be used as measured
project results.

Retain raw data, code versions, calibration details and test sheets in appendices
or permitted digital attachments. Document individual contributions, guide
feedback and action closure. Discuss relevant safety, ethics, environmental
effects, sustainability, cost and standards using evidence. An SDG label alone
is insufficient.

\subsection*{11 Before guide approval}
\begin{itemize}
\item Replace every placeholder and check title, names, roll numbers and degree wording.
\item Use the approved emblem and verify Times New Roman is actually embedded.
\item Check the 500-word abstract, additional report layout and keywords. Keep the
electronic abstract consistent with the report.
\item Verify the correspondence between objectives, requirements, tests and conclusions.
\item Check contents, figure/table lists, captions, equations and bibliography.
\item Obtain genuine signatures; do not simulate approvals in the template.
\item Resolve missing images, undefined citations, text overflow and font substitution.
\end{itemize}

\subsection*{12 Printing and binding}
Print BE reports double-sided on A4 white bond paper of at least 70 gsm and
whiteness of 95\% or above. Print at actual size, not Fit to page. Final hard
binding after viva is black with gold embossed cover and spine lettering. The
spine includes degree, candidate name and year. Ask for two white sheets at the
beginning and end, in addition to normal binding sheets. Confirm team-name spine
arrangement and final copy requirements with the department; no copy count is
invented here.

\Page{Source Notes and Formatting Decisions}
\subsection*{13 Source requirements and recommended content}
Source: \textit{Guidelines for Preparation of Thesis, Dissertations and Project
Reports}, user-provided 16-page PDF. Content structure is on PDF pages 2--4;
formatting on pages 5--7; auxiliary pages and binding on pages 7--10; specimens
A--F on pages 11--16. Page numbers here count from the cover.

Introduction first, Literature Review normally second, Results and Discussions
penultimate and Conclusions last are source rules. The exact titles and number
of intervening chapters depend on the investigation. The eight-chapter plan is a
recommended BE structure, not a prescribed university index.

\subsection*{14 Interpretations to confirm with the guide before binding}
\begin{itemize}
\item \Key{Chapter top spacing}: the source says an additional top margin of
about 75 mm. The template applies 75 mm above the numbered chapter heading,
relative to the normal 30 mm top margin, with 12 mm gaps between number, title
and body. If the department treats 75 mm as a total top distance, adjust the
central chapter spacing in settings.tex after confirmation.
\item \Key{Paragraph separation}: the source requests about 2.5-line spacing
between paragraphs and normally a 12 mm first-line indent. The template uses a
12 mm indent and an interpreted paragraph baseline gap of 2.5 ordinary baselines.
Confirm this approximate spacing before final printing.
\item \Key{Duplex margins}: the source gives physical left 30 mm and right
20 mm. The template preserves these on both sides. Confirm any mirrored
inner/outer binding-margin variation with the guide.
\item \Key{Headers}: the source allows no header with a 30 mm top margin. This
option is used, retaining the 160 by 245 mm writing area and centred footer.
\item \Key{Approval sequence}: the source has inconsistent numbering words
around iii. A continuous lower-case Roman sequence is used. Examiner approval
is disabled until the final post-viva version.
\end{itemize}

\subsection*{15 Adapting the package responsibly}
The certificate and declaration specimens are adapted to a four-student team.
Confirm final wording with the guide. For a different team size, edit all name
lists and signature macros, not just the title sheet. A logo area and blank
signature lines are placeholders, not authorised emblems or approvals.

The preview is an instructional template, not a completed thesis. It contains no
fabricated results or real project bibliography. This handout is an instruction
sheet; thesis formatting is implemented in main.tex and settings.tex.
\end{document}
